\documentclass[11pt,a4paper]{article}
\usepackage[T1]{fontenc}
\usepackage[utf8]{inputenc}
\usepackage[french]{babel}
\usepackage{lmodern,amsmath,amssymb,booktabs}
\usepackage[margin=25mm]{geometry}
\usepackage{microtype}
\usepackage[hidelinks]{hyperref}
\newcommand{\ii}{\mathrm{i}}
\newcommand{\E}{\mathbb{E}}
\newcommand{\Reel}{\operatorname{Re}}
\newcommand{\Var}{\operatorname{Var}}
\newcommand{\Cov}{\operatorname{Cov}}
\newcommand{\CN}{\mathcal{CN}}
\title{Mesure d'une cible et d'une épaisseur\\par balayage fréquentiel\\
Modèles de signal, méthodes A et B et information de Fisher}
\author{Projet MicroWaveTomography --- banc plexiglas}
\date{Mise à jour : 18 septembre 2026}
\begin{document}
\maketitle

\begin{abstract}
Nous présentons d'abord le modèle d'une cible unique sans épaisseur :
retard, amplitude complexe, réponse de l'antenne et changements du fond.
Les méthodes A (profil IFFT) et B (ajustement fréquentiel) sont écrites
explicitement, avec leurs limites d'identifiabilité.
Pour une plaque diélectrique, on cherche ensuite à estimer l'épaisseur $e$ et un facteur réel
\textbf{d'amplitude de l'écho arrière} $\alpha$, sans dimension, à partir
des valeurs complexes de $S_{11}$ mesurées à plusieurs fréquences.
Le modèle à deux échos conduit à une matrice d'information de Fisher (FIM)
explicite. Nous distinguons le cas d'une calibration connue du cas où le
gain et le retard jusqu'à la plaque sont aussi inconnus. Une extension
inclut les réflexions internes multiples. Aucune précision expérimentale
n'est chiffrée sans mesure du bruit et du couplage.
\end{abstract}

\section{Cible unique sans épaisseur : signal et méthodes A et B}
\label{sec:cible-unique}

Cette section concerne une cible représentée par \textbf{une seule surface
réfléchissante}, par exemple la face visible d'une plaque métallique.
Il n'y a ni épaisseur $e$ à estimer, ni écho de face arrière à imposer.
Une cible géométriquement unique peut néanmoins engendrer plusieurs trajets
avec la pièce ou une réponse fréquentielle complexe. Le modèle à un seul
écho est donc une approximation à vérifier, et non une conséquence automatique
du nombre d'objets.

\subsection{Ce que mesure le LiteVNA : du signal reçu à S11}
\label{sec:s11-acquisition}

\paragraph{Trois indices différents.}
On note $k=0,\ldots,K-1$ l'indice de fréquence du balayage
(\textbf{$K=201$ dans nos mesures}), $r$ l'indice du balayage répété,
et $m$ l'indice du profil en retard obtenu après IFFT. $K$ est le nombre
de fréquences réellement mesurées ; $M=N_{\rm FFT}$ sera le nombre de
points de l'IFFT. Dans les sections ultérieures consacrées à l'épaisseur,
la notation $N$ pour le nombre de fréquences désigne la même quantité que $K$.

\paragraph{Définition physique, avant de parler de Fourier.}
À une fréquence fixée $f_k$, la référence incidente et le signal réfléchi
peuvent se représenter par des sinusoïdes. En notant $u$ le temps physique
pendant la mesure à cette fréquence, leurs amplitudes complexes sont
$a_{1,k}=A_k\exp(\ii\phi_{a,k})$ et
$b_{1,k}=B_k\exp(\ii\phi_{b,k})$. Les signaux correspondants s'écrivent,
avec une même normalisation des ondes,
\begin{equation}
 a_k(u)=\Reel\{a_{1,k}\exp(\ii2\pi f_ku)\},\qquad
 b_k(u)=\Reel\{b_{1,k}\exp(\ii2\pi f_ku)\}.
\end{equation}
Pour une mesure à un port, sans onde incidente venant d'un second port,
le coefficient de réflexion au plan de référence est
\begin{equation}\label{eq:s11-definition}
 \boxed{S_{11}(f_k)=\frac{b_{1,k}}{a_{1,k}}
 =\frac{B_k}{A_k}\exp[\ii(\phi_{b,k}-\phi_{a,k})].}
\end{equation}
C'est \textbf{un scalaire complexe par fréquence}, donc deux nombres réels,
pas seulement un module ou une puissance. La phase est celle de l'onde
réfléchie \emph{par rapport} à la référence incidente. L'appareil effectue
la détection cohérente ; le programme Python ne reçoit pas ici les échantillons
temporels de la sinusoïde RF et n'effectue pas une première FFT pour les analyser.
Les deux sinusoïdes ci-dessus sont une représentation physique, non une
description de l'algorithme de démodulation interne du firmware.

\paragraph{Valeurs effectivement reçues en mode brut USB.}
Avec l'option \texttt{--sol}, le programme demande le mode brut du LiteVNA.
Pour chaque $k$ et chaque balayage $r$, le paquet USB fournit notamment quatre
entiers représentant les composantes réelle et imaginaire des canaux de
référence et de réflexion~\cite{litevnausb}. Le pilote les assemble en
\begin{equation}
 U_{k,r}=I_{{\rm ref},k,r}+\ii Q_{{\rm ref},k,r},\qquad
 V_{k,r}=I_{{\rm refl},k,r}+\ii Q_{{\rm refl},k,r}.
\end{equation}
$I$ et $Q$ sont des valeurs numériques de détection, non des tensions en volts
calibrées directement au raccord. Le calcul de \texttt{litevna.decoder} est
\begin{equation}\label{eq:s11-brut}
 \boxed{\gamma^{\rm brut}_{k,r}=\frac{V_{k,r}}{U_{k,r}}
 =\frac{I_{\rm refl}+\ii Q_{\rm refl}}{I_{\rm ref}+\ii Q_{\rm ref}}.}
\end{equation}
Les indices $(k,r)$ sont omis dans le dernier quotient pour alléger l'écriture.
Explicitement,
\begin{align}
 \Reel\gamma^{\rm brut}
 &=\frac{I_{\rm refl}I_{\rm ref}+Q_{\rm refl}Q_{\rm ref}}
 {I_{\rm ref}^2+Q_{\rm ref}^2},\\
 \operatorname{Im}\gamma^{\rm brut}
 &=\frac{Q_{\rm refl}I_{\rm ref}-I_{\rm refl}Q_{\rm ref}}
 {I_{\rm ref}^2+Q_{\rm ref}^2}.
\end{align}
Par exemple, $U=1000$ et $V=100-100\ii$ donnent
$\gamma^{\rm brut}=0{,}1-0{,}1\ii$, un module $0{,}1414$ et une phase
$-45^\circ$. Ce quotient brut contient encore les erreurs instrumentales.
Le calcul est indépendant à chaque fréquence : il ne mélange pas les
201 fréquences.

\paragraph{Correction SOL fréquence par fréquence.}
Le modèle d'erreur à un port utilisé dans \texttt{commun/sol.py} est
\begin{equation}
 \gamma^{\rm brut}_{k,r}=e_{00,k}+
 \frac{e_{R,k}\,S_{11,k,r}}
 {1-e_{11,k}S_{11,k,r}}.
\end{equation}
Ici $e_{R,k}=(e_{10}e_{01})_k$ est le coefficient complexe de suivi de
réflexion, stocké sous le nom \texttt{e10e01}. $e_{00,k}$ décrit la directivité
résiduelle et $e_{11,k}$ la désadaptation de source. Les trois coefficients sont déterminés par SHORT, OPEN
et LOAD à chaque fréquence. La mesure corrigée est
\begin{equation}\label{eq:s11-sol}
 \boxed{\widehat S_{11,k,r}=
 \frac{\gamma^{\rm brut}_{k,r}-e_{00,k}}
 {e_{R,k}+e_{11,k}(\gamma^{\rm brut}_{k,r}-e_{00,k})}.}
\end{equation}
Dans le reste du document, $S^{(0)}_{k,r}$ et $S^{(1)}_{k,r}$ désignent ces
valeurs \emph{mesurées et corrigées}, respectivement sans et avec cible.
Elles ne sont pas les valeurs théoriques calculées à partir du modèle de cible.

L'autre mode, \texttt{--interne}, demande au firmware sa propre correction.
Dans la convention USB vérifiée par le pilote, la référence vaut alors
$2^{30}+0\ii$ et la valeur corrigée reçue est
$(I_{\rm refl}+\ii Q_{\rm refl})/2^{30}$. On n'applique pas en plus la SOL
logicielle. Les acquisitions analysées avec \texttt{--sol} suivent bien
les étapes~\eqref{eq:s11-brut} puis~\eqref{eq:s11-sol}.

\paragraph{Des 201 valeurs fréquentielles au profil en retard.}
Après correction de chaque balayage, on moyenne les nombres complexes puis
on soustrait le fond :
\begin{equation}\label{eq:s11-vers-y}
 y_k=\frac1{R_1}\sum_{r=1}^{R_1}\widehat S^{(1)}_{11,k,r}
     -\frac1{R_0}\sum_{r=1}^{R_0}\widehat S^{(0)}_{11,k,r},
 \quad k=0,\ldots,K-1.
\end{equation}
Le vecteur $\mathbf y$ contient toujours \textbf{201 valeurs complexes indexées
par fréquence}. Prendre $20\log_{10}|y_k|$ donne un graphique \emph{en fréquence},
mais ne donne pas la distance. A conserve les phases pour effectuer ensuite
une IFFT ; B ajuste directement ces mêmes valeurs complexes.
Ainsi, la chaîne de calcul est
\begin{equation*}
 (I_{\rm ref},Q_{\rm ref},I_{\rm refl},Q_{\rm refl})_{k,r}
 \ \longrightarrow\ \gamma^{\rm brut}_{k,r}
 \ \longrightarrow\ \widehat S_{11,k,r}
 \ \longrightarrow\ y_k
 \ \begin{cases}
   \text{IFFT puis module : méthode A},\\
   \text{ajustement complexe : méthode B}.
 \end{cases}
\end{equation*}

\subsection{Référence de distance et réponse de la cible}

On mesure aux fréquences $f_k=f_{\min}+k\Delta f$, $k=0,\ldots,K-1$,
avec $B=(K-1)\Delta f$. Le VNA mesure des rapports complexes à chaque fréquence ;
il ne déclenche pas un chronomètre pour chaque écho. La transformation
fréquence--temps est effectuée ensuite~\cite{keysight}.

La distance géométrique $d$ est mesurée d'un repère déclaré sur l'antenne
à la face de la cible, le long de l'axe de mesure. Dans l'air, le retard
de propagation aller-retour est
\begin{equation}
 \tau_d=\frac{2d}{c}.
\end{equation}
Après calibration SOL au raccord de l'antenne, les effets du câble et des
raccords inclus dans cette calibration sont normalement corrigés. La réponse
de l'antenne, absente pendant la SOL, reste présente. Dans une approximation
à trajet direct, on écrit la contribution de la cible
\begin{equation}\label{eq:cible-directe}
 T_k(d)=g\,H_k\exp\left(-\ii2\pi f_k\frac{2d}{c}\right),
 \qquad g=\alpha_{\rm c}\exp(\ii\phi),\quad \alpha_{\rm c}\geq0.
\end{equation}
$H_k$ représente la forme fréquentielle du transfert émission--réception,
incluant la réponse de l'antenne et, si elle est connue, celle de la diffusion
par la cible. $g$ est un gain complexe constant sur la bande ajustée ;
son module $\alpha_{\rm c}$ est une amplitude globale d'écho, \emph{pas}
un coefficient de perte par mètre. Sa valeur dépend de la normalisation de
$H_k$. Il ne faut pas confondre $\alpha_{\rm c}$ avec le facteur $\alpha$
de l'écho arrière utilisé plus loin pour une plaque diélectrique.

Une forme fréquentielle indépendante de $d$ est elle-même une hypothèse :
la diffraction, l'angle, la taille de la plaque et le champ proche peuvent
la modifier. Le gain constant peut changer d'une position de cible à l'autre.

Si, sur la bande considérée, on peut factoriser
\begin{equation}
 H_k=\widetilde H_k\exp(-\ii2\pi f_k\tau_{\rm a}),
 \qquad \tau_{\rm app}=\frac{2d}{c}+\tau_{\rm a},
\end{equation}
alors $\tau_{\rm a}$ est le retard effectif de la chaîne antenne en
aller-retour par rapport au repère choisi. La conversion en distance donne
\begin{equation}\label{eq:cible-offset}
 d_{\rm app}=\frac{c\tau_{\rm app}}{2}
 =d+\frac{c\tau_{\rm a}}{2}.
\end{equation}
Cette relation décrit un retard pur. Si $\widetilde H_k$ porte encore une
phase dispersive, ou si d'autres contributions se superposent, la position
du maximum temporel peut aussi dépendre de la bande et du fenêtrage.
L'offset n'est donc pas garanti constant. Par exemple, 20~cm de distance
affichée correspondent à 1,33~ns et 38~cm à 2,54~ns, sans prouver la présence
d'un objet à ces distances. Une extension de port compense un retard déclaré ;
elle n'identifie pas à elle seule la réponse complète de l'antenne~\cite{portextension}.

\subsection{Fond, occultation et différence réellement traitée}

Soient $S^{(0)}_{k,r}$ le balayage $r$ sans cible et $S^{(1)}_{k,r}(d)$
un balayage avec cible. Un modèle additif explicite est
\begin{align}
 S^{(0)}_{k,r} &= A^{(0)}_k+F^{(0)}_k+\varepsilon^{(0)}_{k,r},\\
 S^{(1)}_{k,r}(d) &= A^{(1)}_k(d)+F^{(1)}_k(d)
                    +T_k(d)+M_k(d)+\varepsilon^{(1)}_{k,r}.
\end{align}
$A$ désigne la contribution propre au montage, $F$ celle de l'environnement,
et $M$ les trajets supplémentaires impliquant la cible. Les contributions
propres à l'antenne peuvent changer si la présence de la cible modifie son
couplage. Cette décomposition est un modèle de travail : les termes ne sont
pas séparément observés à partir d'une seule mesure de $S_{11}$.

Pour $R_0$ balayages de fond et $R_1$ balayages avec cible, la donnée utilisée
par les deux méthodes est
\begin{equation}\label{eq:cible-difference}
 \boxed{y_k=\overline S^{(1)}_k-\overline S^{(0)}_k
 =T_k(d)+\underbrace{\Delta A_k+\Delta F_k+M_k(d)}_{b_k(d)}+\eta_k,}
\end{equation}
où $\Delta A=A^{(1)}-A^{(0)}$ et $\Delta F=F^{(1)}-F^{(0)}$.
La soustraction est faite \textbf{en complexe, avant le module et les dB}.
Elle mesure un changement de scène et n'isole $T$ que si $b$ est négligeable
ou correctement modélisé.

\paragraph{Cas simplifié d'un fond masqué.}
Si la présence de la plaque multiplie une contribution de fond par un facteur
$m_k(d)$, alors $F^{(1)}_k=m_k(d)F^{(0)}_k$ et
\begin{equation}
 \Delta F_k=[m_k(d)-1]F^{(0)}_k.
\end{equation}
Dans ce modèle simplifié, $m=1$ signifie inchangé et $m=0$ masqué complètement.
En général le changement peut aussi porter sur la phase et ne se réduit pas
à une atténuation réelle. Pour une occultation complète et $\Delta A=M=0$,
$y=T-F^{(0)}+\eta$ : le fond est présent avec un signe négatif dans le résidu.
Le module de sa transformée conserve une bosse, mais \textbf{le signe négatif
ne change pas son retard}. Il ne suffit donc pas à expliquer une bosse précoce
par un obstacle situé derrière la plaque.

\paragraph{Cas simplifié d'une dérive du montage.}
Si $A^{(1)}_k=A^{(0)}_k\exp(-\ii2\pi f_k\delta\tau)$, alors
\begin{equation}
 \Delta A_k=A^{(0)}_k[\exp(-\ii2\pi f_k\delta\tau)-1]
 \simeq-\ii2\pi f_k\delta\tau A^{(0)}_k,
 \quad |2\pi f_k\delta\tau|\ll1.
\end{equation}
Un petit changement de phase d'une contribution forte peut donc dominer
un écho faible. Cela décrit une possibilité, pas un diagnostic du connecteur
ou de l'oscillateur. Une contribution fixe dans les quatre acquisitions peut
venir du matériel \emph{ou} d'un environnement fixe ; des modules superposés
ne garantissent pas des phases identiques.

\subsection{Modèle statistique de la moyenne complexe}

On écrit
\begin{equation}\label{eq:cible-stat}
 \mathbf y=\boldsymbol\mu(\boldsymbol\theta)+\mathbf b+
 \boldsymbol\eta,\qquad
 \boldsymbol\eta\sim\CN(\mathbf0,\mathbf C).
\end{equation}
Sous l'hypothèse de balayages indépendants et de mesures avec/sans cible
indépendantes, les covariances par balayage $\mathbf C_1$ et $\mathbf C_0$
donnent
\begin{equation}
 \mathbf C=\frac{\mathbf C_1}{R_1}+\frac{\mathbf C_0}{R_0}.
\end{equation}
Dans l'approximation diagonale, $v_k=\sigma_{1,k}^2/R_1+
\sigma_{0,k}^2/R_0$. Une dérive cohérente ou une occultation ne disparaît
pas nécessairement par moyennage et ne doit pas être assimilée sans justification
à ce bruit indépendant. Le réemploi d'un même fond rend aussi les différences
$\mathbf y$ de plusieurs positions corrélées entre elles.

\subsection{Méthode A : reconstruction non paramétrique}

Le symbole $\tau$ ci-dessous est un \textbf{retard testé}, pas le temps $u$
pendant lequel le VNA mesure une sinusoïde. Il n'est pas fourni comme un
échantillon temporel par le VNA. Pour une fenêtre $w_k$ (Hann dans le programme),
on définit mathématiquement un profil en retard
\begin{equation}\label{eq:cible-A}
 h_w(\tau)=\frac1K\sum_{k=0}^{K-1}w_k y_k\exp(\ii2\pi f_k\tau),
 \qquad p_A(\tau)=|h_w(\tau)|.
\end{equation}
Cette somme combine les $K=201$ fréquences pour chaque retard testé. Le
facteur $1/K$ correspond à la normalisation actuelle du code ; le gain
cohérent de la fenêtre n'est pas compensé.

\paragraph{Calcul discret réellement effectué.}
On choisit $M\geq K$ et on complète $w_k y_k$ par des zéros pour $K\leq k<M$.
Notons ce vecteur $q_k$. La transformée inverse discrète est
\begin{equation}
 z_m=\operatorname{IFFT}_M(\mathbf q)_m
 =\frac1M\sum_{k=0}^{M-1}q_k\exp(\ii2\pi km/M),
 \qquad m=0,\ldots,M-1.
\end{equation}
\textbf{$k$ indexe les fréquences d'entrée ; $m$ indexe les retards de sortie.}
Les valeurs $z_m$ sont encore complexes. Leur module fournit le profil,
avec la grille et la normalisation suivantes :
\begin{equation}
 \tau_m=\frac{m}{M\Delta f},\qquad
 d_m=\frac{c\tau_m}{2},\qquad
 p_A(\tau_m)=\frac{M}{K}
 |z_m|.
\end{equation}
Sans bourrage de zéros, $M=K=201$ : les 201 valeurs d'entrée sont des
fréquences et les 201 valeurs de sortie sont des retards. Dans le programme,
$M=16K=3216$ ; seule la première moitié du profil est affichée. Pour un sweep
de 1 à 6~GHz, $\Delta f=25$~MHz. L'espacement des retards vaut alors
12,44~ps (1,86~mm sur l'axe radar), contre 199,0~ps sans bourrage de zéros.
Ces espacements sont ceux de la grille de calcul, pas la résolution physique.

La phase porteuse $\exp(\ii2\pi f_{\min}\tau_m)$ omise par l'IFFT ne change
pas le module. Le code affiche la première moitié de la période reconstruite.
Le bourrage de zéros affine l'échantillonnage du profil, sans augmenter la
résolution physique ; celle-ci est de l'ordre de $1/B$ en temps, avec un
facteur dépendant de la fenêtre et du critère choisi~\cite{keysight}.
La période d'ambiguïté est $1/\Delta f$, soit $c/(2\Delta f)$ en distance.

Dans une région $\mathcal I$ où l'on a identifié la contribution à suivre,
\begin{equation}
 \widehat\tau_A=\underset{\tau\in\mathcal I}{\operatorname{arg\,max}}\,p_A(\tau),
 \qquad \widehat d_{A,\rm app}=c\widehat\tau_A/2.
\end{equation}
Cette règle de suivi n'est pas une preuve que le plus grand pic est la cible.
Pour un seul écho pur $y_k=g\exp(-\ii2\pi f_k\tau_0)$,
\begin{equation}
 h_w(\tau)=g\,W(\tau-\tau_0),\qquad
 W(u)=\frac1K\sum_kw_k\exp(\ii2\pi f_ku).
\end{equation}
Avec la réponse $H_k$ et les termes $b_k$, on obtient une somme de réponses
filtrées par la bande et la fenêtre. Les lobes secondaires et les interférences
peuvent produire plusieurs maxima : \textbf{un maximum n'est pas nécessairement
un trajet distinct}. Un offset empirique peut être retiré de $\widehat d_A$
seulement si sa validité a été vérifiée sur des distances connues.

\subsection{Méthode B : ajustement paramétrique d'une cible}

Le modèle minimal, avec forme fréquentielle connue $\widetilde H_k$, est
\begin{equation}\label{eq:cible-B}
 \mu_k(g,\tau)=g\,s_k(\tau),\qquad
 s_k(\tau)=\widetilde H_k\exp(-\ii2\pi f_k\tau).
\end{equation}
Il suppose que le terme $b_k$ de~\eqref{eq:cible-difference} a été retiré,
est négligeable, ou a été inclus dans un modèle auxiliaire contraint.
Si toute la phase de l'antenne est connue et contenue dans $\widetilde H$,
$\tau$ peut représenter $2d/c$ ; sinon il s'agit d'un retard apparent.
Pour $\mathbf C$ connue et indépendante des paramètres, le maximum de
vraisemblance gaussien minimise
\begin{equation}
 (\widehat g,\widehat\tau)=
 \underset{g\in\mathbb C,\,\tau\in\mathcal T}{\operatorname{arg\,min}}
 (\mathbf y-g\mathbf s(\tau))^H\mathbf C^{-1}
 (\mathbf y-g\mathbf s(\tau)).
\end{equation}
Pour chaque retard, le gain s'élimine analytiquement :
\begin{align}
 \widehat g(\tau)&=\frac{\mathbf s(\tau)^H\mathbf C^{-1}\mathbf y}
 {\mathbf s(\tau)^H\mathbf C^{-1}\mathbf s(\tau)},\\
 \widehat\tau_B&=\boxed{\underset{\tau\in\mathcal T}{\operatorname{arg\,max}}
 \frac{|\mathbf s(\tau)^H\mathbf C^{-1}\mathbf y|^2}
 {\mathbf s(\tau)^H\mathbf C^{-1}\mathbf s(\tau)}.}
\end{align}
On en déduit $\widehat\alpha_{\rm c}=|\widehat g|$,
$\widehat\phi=\arg\widehat g$ et $\widehat d_{B,\rm app}=c\widehat\tau_B/2$.
L'intervalle $\mathcal T$ doit tenir compte des ambiguïtés de la grille et
les minima concurrents doivent être examinés.

\paragraph{Lien précis entre A et B.}
Pour un écho, $\widetilde H_k=1$ et $\mathbf C=v\mathbf I$, le critère B
se réduit à maximiser $|\sum_k y_k\exp(\ii2\pi f_k\tau)|^2$.
Il coïncide avec la recherche du maximum du profil A \emph{sans fenêtre},
à l'interpolation et à l'optimisation continue près. A avec Hann et B à poids
uniformes n'emploient pas la même pondération. B n'est donc pas automatiquement
plus juste ; son intérêt dépend du modèle et des contraintes physiques.

\paragraph{Plusieurs contributions, toujours sans épaisseur.}
Pour $J$ trajets supposés, une extension est
\begin{equation}
 \mu_k=\sum_{j=1}^{J}g_j\widetilde H_{j,k}
               \exp(-\ii2\pi f_k\tau_j).
\end{equation}
Les gains complexes sont obtenus par moindres carrés linéaires pour chaque
ensemble de retards, puis les retards sont optimisés. $J=2$ impose deux
contributions mathématiques, pas deux faces d'une plaque métallique.
Deux retards qui coïncident rendent leurs gains individuels non identifiables.
Le résidu complexe, les solutions aux bornes et le rang de la FIM doivent
être contrôlés avant d'interpréter les paramètres comme des trajets physiques.

\paragraph{Correspondance avec les programmes.}
\texttt{s11\_cible\_A.py} utilise~\eqref{eq:cible-A} avec Hann.
Le programme \texttt{s11\_cible\_B.py}, avec les options
\texttt{--modele echos} et \texttt{--echos 1}, ajuste un écho,
et \texttt{--echos 2} en ajuste deux. Le code actuel suppose
$\widetilde H_{j,k}=1$ : il n'identifie pas la réponse de l'antenne.
Son centrage $f_k-f_c$ au lieu de $f_k$ est absorbé par la phase des gains
libres, sans modifier les retards. Il utilise des poids uniformes par défaut ;
\texttt{--ponderation bruit} emploie $1/v_k$ si ces variances sont disponibles.
Le modèle général avec $\widetilde H$ connue ci-dessus est une formulation
théorique, pas une option de déconvolution déjà présente dans le programme.

\subsection{Identifiabilité et information de Fisher pour une cible}

Si le retard de l'antenne $\tau_{\rm a}$ et la distance $d$ sont tous deux
inconnus, l'écho pur ne dépend que de $2d/c+\tau_{\rm a}$. Pour toute variation
$\delta d$, la transformation
\begin{equation}
 d'=d+\delta d,\qquad \tau_{\rm a}'=\tau_{\rm a}-2\delta d/c
\end{equation}
laisse le signal inchangé. \textbf{Un ajustement sur une seule position ne peut
pas séparer ces deux inconnues.} De même, une réponse complexe arbitraire
$H_k$ libre à chaque fréquence peut absorber n'importe quel retard proposé.
Une caractérisation indépendante de l'antenne ou des distances de référence
est nécessaire pour fixer le zéro géométrique.

Dans le modèle identifiable~\eqref{eq:cible-B}, posons
$\boldsymbol\theta=(\alpha_{\rm c},\phi,\tau)$, avec
$g=\alpha_{\rm c}\exp(\ii\phi)$ et $\alpha_{\rm c}>0$. Les dérivées sont
\begin{equation}
 \partial_{\alpha_{\rm c}}\mu_k=\mu_k/\alpha_{\rm c},\qquad
 \partial_\phi\mu_k=\ii\mu_k,\qquad
 \partial_\tau\mu_k=-\ii2\pi f_k\mu_k.
\end{equation}
Pour un bruit complexe propre de covariance connue,
\begin{equation}
 F_{ab}=2\Reel\!\left[(\partial_a\boldsymbol\mu)^H
                  \mathbf C^{-1}(\partial_b\boldsymbol\mu)\right].
\end{equation}
Si $\mathbf C=\operatorname{diag}(v_k)$, définissons
$\rho_k=|\mu_k|^2/v_k$ et
$\overline f_\rho=\sum_k\rho_k f_k/\sum_k\rho_k$.
Après élimination de la phase inconnue (et de l'amplitude, orthogonale ici),
\begin{equation}
 F_{\tau,\rm eff}=8\pi^2\sum_k\rho_k(f_k-\overline f_\rho)^2,
 \qquad
 \Var(\widehat\tau)\geq F_{\tau,\rm eff}^{-1}.
\end{equation}
Pour une distance avec zéro connu, $d=c\tau/2$, donc
\begin{equation}
 \boxed{\Var(\widehat d)\geq
 \frac{c^2}{32\pi^2\sum_k\rho_k(f_k-\overline f_\rho)^2}.}
\end{equation}
Ce sont des bornes locales pour un estimateur non biaisé sous un modèle
correct. Elles ne bornent ni l'erreur de zéro de l'antenne, ni le biais dû à
un fond mal soustrait. Sans forme fréquentielle calibrée, elles concernent
au mieux le retard apparent du modèle choisi.

\paragraph{Diagnostic sans nouvelle hypothèse d'épaisseur.}
Pour deux positions mesurées avec le même fond numérique,
\begin{equation}
 \mathbf y(d_1)-\mathbf y(d_2)
 =\overline{\mathbf S}^{(1)}(d_1)-\overline{\mathbf S}^{(1)}(d_2).
\end{equation}
Le fond enregistré s'annule exactement, mais les modifications de scène et
les dérives entre positions ne s'annulent pas nécessairement. Comparer cette
différence, plusieurs fenêtres et plusieurs sous-bandes aide à distinguer une
structure commune d'une contribution mobile, sans attribuer automatiquement
chaque bosse à un objet. Une réponse répétable sans cible ne suffit pas à
identifier le composant matériel responsable.

\clearpage
\section{Plaque diélectrique : paramètres et conventions physiques}

Les fréquences du balayage sont
\begin{equation}
 f_k=f_{\min}+k\Delta f,\quad k=0,\ldots,N-1,\qquad
 B=f_{\max}-f_{\min}=(N-1)\Delta f.
\end{equation}
START et STOP sont des fréquences, en hertz. À chaque fréquence, le VNA
mesure une amplitude et une phase, représentées par un nombre complexe.
L'ajustement porte sur ces nombres, avant module, logarithme ou IFFT.

\begin{center}
\begin{tabular}{lll}
\toprule
Symbole & Définition & Unité \\
\midrule
$e>0$ & épaisseur géométrique de la plaque & m \\
$\alpha\geq0$ & facteur d'amplitude de l'écho arrière & sans dimension \\
$n_k$ & indice de phase réel, supposé connu & sans dimension \\
$\beta_k=2\pi f_kn_k/c$ & constante de phase & rad/m \\
$\tau_0$ & retard aller-retour jusqu'à la face avant & s \\
$G_k$ & gain complexe du trajet et de l'instrument & sans dimension \\
\bottomrule
\end{tabular}
\end{center}

\paragraph{Définition précise de $\alpha$.}
Nous définissons $\alpha$ comme le multiplicateur réel de l'amplitude
complexe \emph{nominale} de l'écho arrière, les transmissions et réflexions
aux interfaces étant déjà incluses dans cette amplitude nominale.
Ainsi $\alpha=1$ conserve l'écho nominal ; $\alpha=0{,}5$ divise son
amplitude par deux et sa puissance par quatre. Pour une perte pure,
$0\leq\alpha\leq1$ et la perte aller-retour vaut $-20\log_{10}\alpha$ dB.
On peut autoriser $\alpha>1$ si ce facteur représente aussi un écart au
couplage nominal. Sa phase est fixée par les coefficients physiques connus.

\textbf{$\alpha$ n'est pas un coefficient d'atténuation par mètre.}
Si un coefficient linéique en amplitude $\kappa$ (Np/m) décrit la propagation,
alors, pour un aller-retour dans la plaque,
\begin{equation}\label{eq:conversion}
 \alpha=\exp(-2\kappa e),\qquad
 \kappa=-\frac{\log\alpha}{2e}\quad(\alpha>0).
\end{equation}
Dans le présent modèle statistique, $e$ et $\alpha$ sont estimés comme
\emph{deux paramètres indépendants}. En particulier, on ne dérive pas
$\alpha$ par rapport à $e$ lors du calcul de la FIM en coordonnées $(e,\alpha)$.
Nous supposons $\alpha$ constant sur la bande ; une perte dépendant de
la fréquence nécessite une paramétrisation supplémentaire.

\paragraph{Retard et résolution.}
Avec la convention temporelle $\exp(\ii2\pi ft)$, un trajet aller-retour
interne donne $\exp(-2\ii\beta_ke)$. Pour un indice constant $n$,
\begin{equation}
 \Delta\tau=\frac{2ne}{c}.
\end{equation}
Deux réflecteurs distants de 15~mm dans l'air donnent donc environ 100~ps
d'écart aller-retour (50~ps en trajet simple). Une plaque de 15~mm d'indice
supposé $\sqrt{2{,}6}$ donne environ 161~ps. Une couche d'air au sein du
même air ne réfléchit pas : l'exemple des 100~ps suppose deux réflecteurs.
Pour un indice dispersif, le retard de groupe du trajet interne est
$2e[n+f\,\mathrm{d}n/\mathrm{d}f]/c$.

De 1 à 3,5~GHz, $B=2{,}5$~GHz et $1/B=400$~ps, soit une échelle de
résolution en réflexion de $c/(2B)\simeq60$~mm dans l'air. Le coefficient
exact dépend du fenêtrage et du critère de résolution~\cite{keysight}.
Cette largeur des pics reconstruits n'est pas une borne universelle
sur l'erreur d'estimation de $e$.

\section{Deux approches : non paramétrique (A) et paramétrique (B)}

\paragraph{Méthode A --- non paramétrique.}
On reconstruit un profil temporel par transformée de Fourier inverse des
mesures fréquentielles, sans imposer au préalable une réponse de plaque
décrite par quelques paramètres. On cherche ensuite les pics associés aux
deux faces et l'on convertit leur écart en épaisseur, pour un indice connu.
La séparation des échos dépend de la bande et du fenêtrage. Le qualificatif
non paramétrique concerne la reconstruction du profil : il ne signifie pas
que le traitement est exempt d'hypothèses physiques ou de réglages.

\paragraph{Méthode B --- paramétrique.}
On impose une famille de réponses complexes décrite par un nombre fini de
paramètres, ici $(e,\alpha)$ et, si nécessaire, le gain et le retard externes.
On ajuste ces paramètres directement aux mesures fréquentielles. Deux pics
temporels séparés ne sont pas nécessaires ; l'identifiabilité et la validité
du modèle doivent cependant être établies.

\textbf{La FIM calculée dans ce document concerne la méthode B et le modèle
statistique choisi.} Elle ne constitue pas une évaluation automatique des
performances de la détection de pics de la méthode A. Les deux méthodes
peuvent utiliser les mêmes mesures complexes et se compléter : le profil
temporel aide à examiner la scène, l'ajustement paramétrique estime ses
paramètres sous les hypothèses déclarées.

\section{Méthode B : moyenne physique à deux échos cohérents}\label{sec:modele}

On suppose une plaque plane homogène, une incidence normale, des milieux
non magnétiques et une scène stable pendant le balayage. Les effets de
bord, la diffraction et les changements de trajets parasites ne sont pas
inclus dans ce modèle.

Notons $r_k=r_{12,k}$ la réflexion de la face avant et $s_k=r_{23,k}$
celle de la face arrière vue depuis la plaque. Pour air/plaque/air,
en négligeant l'effet des pertes sur les interfaces,
\begin{equation}
 r_k=\frac{1-n_k}{1+n_k},\qquad s_k=-r_k,\qquad
 b_k=(1-r_k^2)s_k.
\end{equation}
Le facteur $1-r_k^2$ représente les deux transmissions de la face avant.
La moyenne complexe après soustraction du fond est
\begin{equation}\label{eq:mu}
 \boxed{\mu_k(e,\alpha)=P_k\left[r_k+\alpha b_kz_k(e)\right],\qquad
 P_k=G_k\exp(-\ii2\pi f_k\tau_0),\quad
 z_k(e)=\exp(-2\ii\beta_ke).}
\end{equation}
Le premier terme est l'écho de la face avant. Le second est celui de
la face arrière, avec une amplitude multipliée par $\alpha$ et un retard
lié à $e$. Les signes des coefficients de réflexion portent la différence
de phase aux interfaces ; $\alpha$ reste réel positif.

Si l'on préfère définir un rapport d'amplitudes arrière/avant, son module
est $\alpha|b_k/r_k|$ lorsque $r_k\ne0$. Il diffère donc de notre facteur
$\alpha$, qui est référencé à l'écho arrière nominal.

La soustraction d'un fond acquis sans cible ne retire pas l'écho de la
face avant. Elle retire les contributions présentes dans le fond, sous
l'hypothèse que l'insertion de la plaque ne modifie pas ces contributions.
Le modèle à deux échos néglige les allers-retours internes supplémentaires ;
la section~\ref{sec:multiple} donne le modèle complet.

\section{Modèle statistique des données complexes}

On forme le vecteur $\mathbf y=(y_0,\ldots,y_{N-1})^T$ des moyennes
complexes avec cible, moins les moyennes du fond. On suppose
\begin{equation}\label{eq:stat}
 \boxed{\mathbf y=\boldsymbol\mu(\boldsymbol\theta)+\boldsymbol\varepsilon,
 \quad \boldsymbol\theta=(e,\alpha)^T,\quad
 \boldsymbol\varepsilon\sim\CN(\mathbf0,\mathbf C).}
\end{equation}
Le bruit est gaussien complexe propre :
$\E[\boldsymbol\varepsilon\boldsymbol\varepsilon^H]=\mathbf C$ et
$\E[\boldsymbol\varepsilon\boldsymbol\varepsilon^T]=0$.
La covariance $\mathbf C$ est supposée connue, définie positive et
indépendante des paramètres. Pour des fréquences indépendantes,
\begin{equation}\label{eq:variance}
 \mathbf C=\operatorname{diag}(v_k),\qquad v_k=\E|\varepsilon_k|^2,
 \quad \Var(\Reel\varepsilon_k)=\Var(\operatorname{Im}\varepsilon_k)=v_k/2.
\end{equation}
$v_k$ désigne la variance complexe totale, pas celle d'une seule quadrature.
Cette convention explique le facteur 2 dans la FIM.

\paragraph{Bruit de la moyenne et du fond.}
Pour $M_c$ sweeps cible et $M_f$ sweeps fond indépendants, avec covariances
par sweep $\boldsymbol\Sigma_c$ et $\boldsymbol\Sigma_f$,
\begin{equation}\label{eq:covfond}
 \mathbf C=\frac{\boldsymbol\Sigma_c}{M_c}
          +\frac{\boldsymbol\Sigma_f}{M_f}.
\end{equation}
Augmenter seulement $M_c$ ne supprime pas le bruit du fond. Réutiliser un
même fond dans plusieurs mesures les corrèle entre elles. Une calibration
commune incertaine peut également créer des corrélations. Si les moyennes
cible et fond sont corrélées, on soustrait de~\eqref{eq:covfond} leur
covariance croisée et sa transposée conjuguée. Une dérive déterministe
ne devient pas un bruit indépendant par simple moyennage.

La densité et la log-vraisemblance sont
\begin{align}
 p(\mathbf y;\boldsymbol\theta)&=
 \frac{\exp[-(\mathbf y-\boldsymbol\mu)^H\mathbf C^{-1}
 (\mathbf y-\boldsymbol\mu)]}{\pi^N\det\mathbf C},\\
 \ell(\boldsymbol\theta)&=-(\mathbf y-\boldsymbol\mu)^H\mathbf C^{-1}
 (\mathbf y-\boldsymbol\mu)+\text{constante}.
\end{align}
Le maximum de vraisemblance minimise ce résidu quadratique pondéré
sur un domaine physique en $(e,\alpha)$.

\section{Construction de la matrice d'information de Fisher}

Pour des paramètres réels, posons
$\mathbf d_i=\partial\boldsymbol\mu/\partial\theta_i$.
Le score s'écrit
\begin{equation}
 \partial_{\theta_i}\ell
 =2\Reel\{\mathbf d_i^H\mathbf C^{-1}(\mathbf y-\boldsymbol\mu)\}.
\end{equation}
L'espérance du produit de deux scores donne la formule gaussienne complexe
de Slepian--Bangs~\cite{abeida} :
\begin{equation}\label{eq:fimg}
 \boxed{F_{ij}=2\Reel\{\mathbf d_i^H\mathbf C^{-1}\mathbf d_j\}.}
\end{equation}
Pour une covariance diagonale,
\begin{equation}\label{eq:fim2}
 \boxed{\mathbf F(e,\alpha)=2\sum_k\frac{1}{v_k}
 \begin{pmatrix}
 |\partial_e\mu_k|^2 & \Reel\{(\partial_e\mu_k)^*\partial_\alpha\mu_k\}\\
 \Reel\{(\partial_e\mu_k)^*\partial_\alpha\mu_k\} & |\partial_\alpha\mu_k|^2
 \end{pmatrix}.}
\end{equation}
Si $\mathbf C$ dépend des paramètres, il faut ajouter
$\operatorname{tr}(\mathbf C^{-1}\mathbf C_{,i}\mathbf C^{-1}\mathbf C_{,j})$
à~\eqref{eq:fimg}. Une covariance estimée puis tenue fixe conduit à une
borne conditionnelle approchée.

\section{FIM explicite pour \texorpdfstring{$(e,\alpha)$}{(e, alpha)} : calibration connue}\label{sec:fermee}

Dans cette section, $r_k$, $b_k$, $n_k$, $G_k$ et $\tau_0$ sont connus
et indépendants de $(e,\alpha)$. Définissons l'écho arrière nominal
\begin{equation}
 u_k=P_kb_kz_k(e).
\end{equation}
Les deux dérivées de~\eqref{eq:mu} sont
\begin{equation}\label{eq:deriv}
 \boxed{\partial_e\mu_k=-2\ii\beta_k\alpha u_k,\qquad
        \partial_\alpha\mu_k=u_k.}
\end{equation}
L'épaisseur modifie la phase de l'écho arrière ; $\alpha$ modifie son
amplitude. Leur produit scalaire réel à chaque fréquence est nul :
\begin{equation}
 \Reel\{(\partial_e\mu_k)^*\partial_\alpha\mu_k\}
 =\Reel\{2\ii\beta_k\alpha|u_k|^2\}=0.
\end{equation}
Pour un bruit indépendant entre fréquences, posons
\begin{equation}
 W_0=\sum_k\frac{|u_k|^2}{v_k},\qquad
 W_2=\sum_k\frac{\beta_k^2|u_k|^2}{v_k}.
\end{equation}
On obtient la FIM demandée :
\begin{equation}\label{eq:fimfermee}
 \boxed{\mathbf F(e,\alpha)=
 \begin{pmatrix}
 8\alpha^2W_2 & 0\\
 0 & 2W_0
 \end{pmatrix}.}
\end{equation}
Cette diagonale exprime une orthogonalité locale des sensibilités, sous
ces hypothèses. Elle ne signifie pas que les estimateurs sont toujours
indépendants, ni que $e$ et $\alpha$ resteront découplés quand la calibration
sera inconnue. Avec une covariance fréquentielle générale, on utilise
\eqref{eq:fimg} ; le terme croisé ne s'annule plus nécessairement.

\subsection{Bornes de Cramér--Rao}
Pour un estimateur non biaisé, un modèle régulier et un point intérieur,
\begin{equation}\label{eq:crb}
 \Cov(\widehat{\boldsymbol\theta})\succeq\mathbf F^{-1},\qquad
 \boxed{\Var(\hat e)\geq\frac{1}{8\alpha^2W_2},\quad
        \Var(\hat\alpha)\geq\frac{1}{2W_0}.}
\end{equation}
Le symbole $\succeq$ signifie que la différence est semi-définie positive.
Les unités de $F_{ee}$ et de $F_{\alpha\alpha}$ sont respectivement
$\mathrm{m}^{-2}$ et 1. L'écart-type minimal est la racine carrée de la borne
sur la variance.

Lorsque $\alpha\to0$, l'information sur l'épaisseur disparaît : sans écho
arrière, on ne peut plus mesurer son retard. À $\alpha=0$ exactement, la
FIM est singulière et les bornes régulières ne s'appliquent pas à la paire.
Si l'on impose $\alpha\leq1$, le point $\alpha=1$ est aussi une frontière,
où les bornes du problème contraint demandent un traitement spécifique.

Une FIM inversible caractérise une identifiabilité locale régulière, pas
l'unicité globale. Même une mesure à une fréquence peut avoir une FIM locale
inversible dans ce cas calibré, tout en présentant une ambiguïté d'épaisseur
par multiples de $c/(2nf)$ pour un indice constant. Des minima concurrents,
le bruit élevé ou une erreur de modèle peuvent empêcher d'atteindre les bornes.

\subsection{Expression à l'aide du rapport signal sur bruit}
Posons $\rho_k=|\alpha u_k|^2/v_k$, le rapport signal sur bruit complexe
\emph{de l'écho arrière}, et $\rho_{\rm tot}=\sum_k\rho_k$. Alors
\begin{equation}
 F_{ee}=8\sum_k\rho_k\beta_k^2,\qquad
 F_{\alpha\alpha}=\frac{2\rho_{\rm tot}}{\alpha^2}.
\end{equation}
La présence de la fréquence absolue dans $\beta_k^2$ repose ici sur une
référence de phase connue. La section suivante montre ce qui change lorsque
la phase est inconnue.

\section{Paramètres auxiliaires : gain, phase et retard}

Une SOL au bout du câble ne détermine pas à elle seule le couplage complexe
antenne--plaque. Supposons $G_k=g\,\overline G_k$, avec une forme fréquentielle
$\overline G_k$ connue, mais un gain $g=g_R+\ii g_I$ et un retard $\tau_0$
inconnus. Le vecteur complet est
\begin{equation}
 \boldsymbol\vartheta=(e,\alpha,g_R,g_I,\tau_0)^T.
\end{equation}
Avec $a_k=\overline G_k\exp(-\ii2\pi f_k\tau_0)
[r_k+\alpha b_kz_k]$, les dérivées supplémentaires sont
\begin{equation}
 \partial_{g_R}\mu_k=a_k,\qquad
 \partial_{g_I}\mu_k=\ii a_k,\qquad
 \partial_{\tau_0}\mu_k=-\ii2\pi f_k\mu_k.
\end{equation}
On construit la FIM complète par~\eqref{eq:fimg}. En notant
$\boldsymbol\eta=(g_R,g_I,\tau_0)$ les paramètres auxiliaires,
\begin{equation}
 \mathbf F_{\text{complète}}=
 \begin{pmatrix}\mathbf F_{\theta\theta}&\mathbf F_{\theta\eta}\\
 \mathbf F_{\eta\theta}&\mathbf F_{\eta\eta}\end{pmatrix}.
\end{equation}
Si $\mathbf F_{\eta\eta}$ est inversible, l'information effective est
le complément de Schur :
\begin{equation}\label{eq:schur}
 \boxed{\mathbf F_{\rm eff}=\mathbf F_{\theta\theta}
 -\mathbf F_{\theta\eta}\mathbf F_{\eta\eta}^{-1}\mathbf F_{\eta\theta}.}
\end{equation}
La borne pertinente est $\mathbf F_{\rm eff}^{-1}$, si elle existe.
Le bloc $\mathbf F_{\theta\theta}^{-1}$ supposerait les paramètres auxiliaires
connus. Si un bloc est singulier, il faut identifier les directions non
observables ; une pseudo-inverse n'autorise pas à annoncer une précision
finie dans une direction indéterminée.

\paragraph{Ambiguïté exacte avec un seul écho.}
Si l'on ne dispose que de l'écho arrière, avec $n$ constant,
\begin{equation}
 \mu_k=g\alpha\overline G_k b_k
 \exp[-\ii2\pi f_k(\tau_0+2ne/c)].
\end{equation}
Un gain libre absorbe $\alpha$, et un retard libre absorbe $e$.
Pour $\alpha,\alpha'>0$, prendre $g'=g\alpha/\alpha'$ et
$\tau_0'=\tau_0+2n(e-e')/c$ laisse toutes les moyennes inchangées.
Les deux paramètres sont donc non identifiables dans ce modèle.
La contribution de face avant liée au même gain, ou des contraintes de
calibration indépendantes, peuvent apporter l'information manquante.
Il faut vérifier le rang au point étudié. Un gain complexe arbitraire
à chaque fréquence peut absorber toute la réponse aux fréquences où
celle-ci ne s'annule pas.

\subsection{Phase inconnue : la largeur de bande apparaît explicitement}
Pour isoler ce mécanisme, considérons un seul écho de la forme
\begin{equation}
 \mu_k=\alpha A_k\exp(\ii\phi-2\ii\beta_ke),
\end{equation}
avec $A_k$ connu (incluant un retard externe connu), et $\phi$ inconnu.
Posons $w_k=|A_k|^2/v_k$, $W_0=\sum_kw_k$,
$W_1=\sum_kw_k\beta_k$ et $W_2=\sum_kw_k\beta_k^2$.
Dans l'ordre $(e,\alpha,\phi)$, la FIM est
\begin{equation}
 \mathbf F=
 \begin{pmatrix}
 8\alpha^2W_2 & 0 & -4\alpha^2W_1\\
 0 & 2W_0 & 0\\
 -4\alpha^2W_1 & 0 & 2\alpha^2W_0
 \end{pmatrix}.
\end{equation}
Après élimination de $\phi$,
\begin{equation}\label{eq:phase}
 \boxed{\mathbf F_{\rm eff}=
 \begin{pmatrix}
 8\alpha^2\sum_kw_k(\beta_k-\overline\beta)^2 & 0\\
 0 & 2W_0
 \end{pmatrix},\qquad\overline\beta=W_1/W_0.}
\end{equation}
La fréquence absolue est remplacée par l'étalement des fréquences : une
phase constante inconnue absorbe la phase moyenne, mais pas sa pente.
Pour $n$ constant, des poids égaux et $N$ points équidistants incluant les
bornes du balayage,
\begin{equation}
 \sum_kw_k(\beta_k-\overline\beta)^2
 =W_0\left(\frac{2\pi n}{c}\right)^2
 \frac{B^2(N+1)}{12(N-1)}.
\end{equation}
À grand $N$, avec $\rho_{\rm tot}=\alpha^2W_0$,
\begin{equation}
 \operatorname{sd}(\hat e)\ \gtrsim\
 \frac{\sqrt{3/2}\,c}{2\pi nB\sqrt{\rho_{\rm tot}}}.
\end{equation}
Il s'agit d'une approximation de la borne locale, pas d'une garantie sur
l'erreur réelle. Elle montre le rôle conjoint de la bande et du bruit :
une précision inférieure à la largeur des pics est possible dans un modèle
identifiable. Pour la plaque complète et ses paramètres auxiliaires, on
recalcule~\eqref{eq:schur} au lieu de transposer cette formule particulière.

\section{Extension aux réflexions internes multiples}\label{sec:multiple}

Si $\alpha$ représente la perte d'amplitude de propagation sur un
aller-retour, chaque aller-retour supplémentaire introduit le même facteur.
On pose alors
\begin{equation}
 q_k=\alpha z_k(e),\qquad
 \Gamma_k=\frac{r_k+s_kq_k}{1+r_ks_kq_k},\qquad
 \mu_k=P_k\Gamma_k.
\end{equation}
Le développement
$\Gamma_k=r_k+(1-r_k^2)s_kq_k+O(q_k^2)$ retrouve le modèle à deux échos.
L'approximation est justifiée si $|r_ks_kq_k|\ll1$.
Si $\alpha$ représente seulement une correction empirique du premier
écho arrière, sa réutilisation à chaque aller-retour est une hypothèse
supplémentaire, à justifier.

Pour $r_k$ et $s_k$ connus et indépendants de $(e,\alpha)$, définissons
\begin{equation}
 h_k=P_k\frac{s_k(1-r_k^2)z_k}{(1+r_ks_k\alpha z_k)^2}.
\end{equation}
Alors
\begin{equation}
 \partial_e\mu_k=-2\ii\beta_k\alpha h_k,\qquad
 \partial_\alpha\mu_k=h_k.
\end{equation}
La FIM calibrée reste de la forme~\eqref{eq:fimfermee}, en remplaçant
$u_k$ par $h_k$ dans $W_0$ et $W_2$. Ces poids dépendent maintenant de
$(e,\alpha)$ à travers les interférences multiples.

\paragraph{Limite de l'hypothèse sur les interfaces.}
Un matériau dissipatif possède un indice complexe, donc des coefficients
de Fresnel dépendant des pertes. Pour inclure cet effet, on peut utiliser
$\kappa=-\log\alpha/(2e)$ et
$\widetilde n_{2,k}=n_k-\ii c\kappa/(2\pi f_k)$ dans les coefficients
$r_k$ et $s_k$. Ceux-ci dépendent alors de $e$ et de $\alpha$ : leurs
\emph{dérivées totales} doivent être incluses dans~\eqref{eq:fim2}.
En particulier, pour $t\in\{e,\alpha\}$ et $P_k$ connu,
\begin{equation}
 \partial_t\mu_k=\frac{P_k}{(1+r_ks_kq_k)^2}
 \left[(1-s_k^2q_k^2)\partial_t r_k
 +q_k(1-r_k^2)\partial_t s_k
 +s_k(1-r_k^2)\partial_t q_k\right].
\end{equation}
La diagonale~\eqref{eq:fimfermee} n'est alors plus garantie. La constance de
$n_k$ et de $\alpha$ sur une bande est une hypothèse effective ; une loi
matérielle dispersive doit être explicitée si elle est nécessaire.

\section{Application aux données du projet}

\paragraph{Programmes A et B.}
Le programme \texttt{s11\_cible\_A.py} conserve le profil IFFT avec fenêtre
de Hann. Le programme \texttt{s11\_cible\_B.py} ajuste directement les données
complexes, y compris celles d'une acquisition existante (option
\texttt{--mesure}). La SOL et le fond sont communs aux deux méthodes.
L'option \texttt{--modele plaque --indice n} utilise le modèle à deux faces
ci-dessus avec un gain complexe constant inconnu et un retard externe libre ;
elle estime $e$ et $\alpha$ dans des intervalles déclarés.
Le mode par défaut \texttt{--modele echos} ajuste une somme d'échos à retards
et gains complexes libres, adaptée au diagnostic avec une plaque métallique ;
il n'estime pas une épaisseur diélectrique.
La réponse fréquentielle de l'antenne n'est pas déconvoluée. Les distances
restent apparentes et la qualité des résidus doit être examinée avant toute
interprétation des paramètres ou de la FIM. Voir \texttt{plexiglas/README.md}
pour les commandes, hypothèses de bruit et fichiers produits.


\begin{enumerate}
\item Utiliser \texttt{freqs} et \texttt{residu} de \texttt{cible.npz},
produit par \texttt{s11\_cible\_A.py} ou \texttt{s11\_cible\_B.py}. Le résidu est la différence complexe
entre moyenne cible et moyenne du fond.
\item Estimer la covariance au niveau de cette moyenne. Dans
\texttt{fond.npz}, \texttt{sigma\_complexe} est un écart-type complexe
\emph{par sweep} : son carré doit être divisé par le nombre de sweeps du
fond avant d'être ajouté à la variance de la moyenne cible.
\item Vérifier la stabilité, les corrélations fréquentielles et l'hypothèse
de bruit propre. Si cette dernière échoue, utiliser le vecteur réel empilant
parties réelle et imaginaire, avec sa covariance complète.
\item Préciser les grandeurs calibrées : indice, interfaces, réponse
fréquentielle de l'antenne, gain et retard jusqu'à la face avant.
Un indice inconnu doit être ajouté aux paramètres : la phase de propagation
porte d'abord sur le produit $ne$.
\item Ajuster le modèle complexe sur un domaine physique en recherchant
les minima concurrents ; calculer les dérivées et la FIM au point ajusté.
Éliminer les paramètres auxiliaires avant d'inverser.
\item Examiner le rang après mise à l'échelle des paramètres (par exemple
$e$ en mm), les corrélations et les résidus. Comparer les bornes à la
dispersion de répétitions indépendantes.
\end{enumerate}

\paragraph{Durée des sinusoïdes et information.}
Le temps d'établissement et d'intégration intervient dans le bruit $v_k$
et dans la validité de la mesure en régime établi. Il ne se déduit pas de
START et STOP. Une intégration plus longue peut diminuer les variances,
sans réduire directement la largeur des pics à bande fixée. Augmenter
le nombre de points à durée totale fixée ne garantit pas plus d'information,
car les $v_k$ peuvent augmenter. Sur une grille régulière, $1/\Delta f$
est la période du profil temporel reconstruit, distincte de sa résolution
$\sim1/B$.

\paragraph{Conclusion.}
Pour une calibration connue et un bruit indépendant, l'épaisseur agit
sur la phase et le facteur $\alpha$ sur l'amplitude de l'écho arrière :
la FIM est la matrice diagonale~\eqref{eq:fimfermee}. Pour une mesure réelle,
la question décisive est l'information restant après prise en compte du
gain, du retard et des autres inconnues, donnée par~\eqref{eq:schur}.
La résolution de 400~ps entre 1 et 3,5~GHz ne suffit donc pas à conclure
sur la précision de l'épaisseur : il faut le bruit, le modèle et sa FIM effective.

\begin{thebibliography}{9}
\bibitem{abeida}
H. Abeida et J.-P. Delmas,
\emph{Slepian--Bangs formulas for parameterized density generator of
elliptically symmetric distributions}, Signal Processing 205 (2023), 108886.
\href{https://doi.org/10.1016/j.sigpro.2022.108886}{doi:10.1016/j.sigpro.2022.108886}.
Référence pour la formule gaussienne générale et ses extensions ; les
calculs du modèle de plaque sont développés dans le présent document.
\bibitem{keysight}
Keysight, \emph{Time Domain}, documentation E5080A,
\href{https://helpfiles.keysight.com/csg/e5080a/time/timedomain.htm}
{transformation temporelle, résolution et étendue de mesure}.
\bibitem{portextension}
Keysight, \emph{Port Extensions}, documentation VNA,
\href{https://helpfiles.keysight.com/csg/pxivna/S3_Cals/Port_Extensions.htm}
{déplacement du plan de référence après calibration}.
\bibitem{litevnausb}
\emph{LiteVNA User Guide}, annexe « USB data interface »,
\href{https://nanovna.com/wp-content/uploads/2021/11/LiteVNA_User-Guide.pdf}
{description du protocole USB}. Les expressions de décodage ci-dessus
correspondent au pilote \texttt{plexiglas/progs/litevna.py} utilisé dans le projet.
\end{thebibliography}
\end{document}
